\section{Conclusion}
\label{sec:conclusion}

We implemented and evaluated a prior-aware multimodal chest X-ray classification model that combines radiographs, clinical indication text, vital signs, and prior-study context. The model was evaluated on a filtered CXR-LT-2023-aligned test split.

The proposed model achieved 0.5662 mAP and 0.9026 mAUROC, outperforming the CaMCheX (Image-only) and no-prior multimodal baselines under the same controlled setting. The largest gains were observed for medium and tail labels, supporting the value of clinical context and prior-study information for long-tailed chest X-ray classification. Compared with external CXR-LT-2023 references, the proposed model exceeds the published number of the top CXR-LT-2023 challenge solution CheXFusion and approaches the published CaMCheX result. We stress that these external comparisons are contextual rather than controlled: the published methods were trained on the full CXR-LT-2023 data at higher resolution and evaluated on the official test set, whereas our results come from a smaller filtered cohort at $512 \times 512$ resolution with a local patient-level test split.

The strongest claims of this thesis are therefore the internal, controlled ones. Within a fixed cohort, resolution, and training budget, adding clinical text, vital signs, and prior-study context to an image-only baseline raises mAP from 0.3246 to 0.5662. The modality ablations show that clinical indication text and prior-study context drive this gain, while structured vital signs contribute little on their own. The main open risks, discussed in Subsection~\ref{subsec:discussion_limitations}, are the reduced training scale, potential shortcut reliance on prior context, and the absence of external multi-center validation. Overall, the results suggest that prior-aware multimodal fusion is a practical and effective direction for clinically contextualized chest X-ray classification, with the above caveats defining the scope of that claim.

\subsection{Future Work}
\label{subsec:conclusion_future}

Future research should focus on:

\textbf{Full-Scale Training}: With more compute, training at larger scale with higher-resolution inputs would allow the model to capture fine-grained visual findings and benefit from stronger data augmentation.

\textbf{End-to-End Multimodal Fine-Tuning}: Fine-tuning the clinical text and image encoders jointly would allow tighter semantic alignment between the modalities.

\textbf{High-Resolution and Localized Attention}: Implementing localized attention or patch-based processing would allow the model to process high-resolution radiographs without memory overflow, preserving tiny pathological details.

\textbf{Advanced Temporal Architectures}: Replacing simple cross-attention with recurrent or state-space models (like Mamba \cite{gu2024mamba}) could allow the network to fuse sequences of multiple prior studies rather than a single previous check.

\textbf{Multi-Center Validation}: Evaluating the prior-aware model on external datasets (such as CheXpert or PadChest \cite{bustos2020padchest}) is crucial to verify generalizability across different clinical settings.

\textbf{Incorporation of Additional Physiological Modalities}: Extending the multimodal framework to include other clinical modalities could capture complementary patient states. Specifically, we cite MIMIC-IV-ECG \cite{gow2023mimiciv_ecg} for diagnostic 12-lead electrocardiograms and the MIMIC-IV Waveform Database \cite{moody2022mimiciv_waveform} for high-resolution ICU bedside-monitor waveforms. Both resources are distributed through PhysioNet \cite{goldberger2000physionet}, so we also cite the standard PhysioNet reference.

